\chapter{Introdução}

O presente trabalho se trata do desenvolvimento de uma ferramenta baseada um métodos numéricos, seu uso é variado mas inicialmente é voltado para a aproximação de soluções de problemas multi-física, isto é, mecânica linear e não-linear, difusão de Fick, propagação de ondas, entre outros que ainda não foram implementados. A biblioteca  é o resultado de um trabalho e minha participação foi mais voltada para módulos de mecânica, mas muitos módulos são usados por todos os tipos de problema, então meu trabalho aborde pontos mais e não só o meu assunto foco.
Os módulos de solução são baseadas em métodos discretos (Finite Element Method - FEM) e semi-analíticos (Semi-Analytical Finite Elements - SAFE) presentes nas diversas bibliotecas escritas em MATLAB que juntas formam uma poderosa ferramenta de soluções numéricas. Essas bibliotecas estão divididas em três grandes grupos: Pre-Processamento, Processamento e Pós-Processamento. 
 
A implementação das rotinas de solução de mecânica linear e não-linear.